% GENERATED FILE. Edit canonical lessons, then run npm run generate:books.
% canonical-source-hash: fnv1a64:8d3ef11ba4b94646
% canonical-lessons: MR-C73-ajari, MR-C73-vedana, MR-C73-tahan, MR-C73-bhuk, MR-C73-anandi

\chapter{Ill, Pain, Thirst, Hunger, Happy}
\label{ch:mr-ill-pain-thirst-hunger-happy}
\hlchaptermodality{73}

\begin{chapteropening}
\textbf{By the end of this chapter:} \emph{I can say ill, sick, pain, thirst, hunger and happy.}

Complete the last lesson of chapter 73: I can say ill, sick, pain, thirst, hunger and happy.
\end{chapteropening}

\section[\texorpdfstring{\mr{आजारी} (ājārī)}{आजारी (ājārī)}]{\mr{आजारी} (ājārī) — ill, sick}

\label{lesson:MR-C73-ajari}



\begin{quote}
Before the new one: say the Marathi for a fever, then the Marathi for a clinic.
\end{quote}

\subsection*{You'll want to know: \mr{आजारी}}
\textbf{\mr{आजारी}} — \emph{ājārī} — \textquotedblleft{}ill, sick\textquotedblright{}.

\begin{cousinweb}
\textbf{\mr{आजारी}} (\emph{ājārī}) is a loan, usually traced to Persian \emph{āzār}, \textquotedblleft{}harm, pain, illness\textquotedblright{}. Marathi adds \emph{-ī}: \textquotedblleft{}having an illness\textquotedblright{}.
\end{cousinweb}

\begin{grammarlens}[title={What you've built}]
One more word for this chapter.
\end{grammarlens}

\subsection*{Your turn}
\begin{itemize}
\raggedright
  \item \emph{Say it:} \emph{ājārī}
  \item \emph{Say it:} \emph{ājārī}, once more
  \item \emph{Say it:} say \emph{ājārī}
  \item \emph{Recall:} say the Marathi for a fever, then the Marathi for a clinic, then say \emph{ājārī} again
\end{itemize}

\begin{tcolorbox}[breakable,colback=teal!4,colframe=teal!35!black,title={Before you move on}]
What does \mr{आजारी} mean? (\textquotedblleft{}Ill, sick\textquotedblright{}.) Say it once more.
\end{tcolorbox}
\section[\texorpdfstring{\mr{वेदना} (vedanā)}{वेदना (vedanā)}]{\mr{वेदना} (vedanā) — pain}

\label{lesson:MR-C73-vedana}



\begin{quote}
Before the new one: say the Marathi for a clinic, then the Marathi for ill.
\end{quote}

\subsection*{You'll want to know: \mr{वेदना}}
\textbf{\mr{वेदना}} — \emph{vedanā} — \textquotedblleft{}pain\textquotedblright{}.

\begin{cousinweb}
\textbf{\mr{वेदना}} (\emph{vedanā}) is Sanskrit \emph{vedanā}, \textquotedblleft{}feeling; pain\textquotedblright{}, taken over whole as a learned loan, a tatsama (CDIAL 12106). CDIAL lists the inherited Marathi form as \emph{veṇ}, \textquotedblleft{}pain, esp. of birth\textquotedblright{}.
\end{cousinweb}

\begin{grammarlens}[title={What you've built}]
One more word for this chapter.
\end{grammarlens}

\subsection*{Your turn}
\begin{itemize}
\raggedright
  \item \emph{Say it:} \emph{vedanā}
  \item \emph{Say it:} \emph{vedanā}, once more
  \item \emph{Say it:} say \emph{vedanā}
  \item \emph{Recall:} say the Marathi for a clinic, then the Marathi for ill, then say \emph{vedanā} again
\end{itemize}

\begin{tcolorbox}[breakable,colback=teal!4,colframe=teal!35!black,title={Before you move on}]
What does \mr{वेदना} mean? (\textquotedblleft{}Pain\textquotedblright{}.) Say it once more.
\end{tcolorbox}
\section[\texorpdfstring{\mr{तहान} (tahān)}{तहान (tahān)}]{\mr{तहान} (tahān) — thirst}

\label{lesson:MR-C73-tahan}



\begin{quote}
Before the new one: say the Marathi for ill, then the Marathi for pain.
\end{quote}

\subsection*{You'll want to know: \mr{तहान}}
\textbf{\mr{तहान}} — \emph{tahān} — \textquotedblleft{}thirst\textquotedblright{}.

\begin{cousinweb}
\textbf{\mr{तहान}} (\emph{tahān}) is a tadbhava, a word inherited from Sanskrit and reshaped by sound change. It comes from \emph{tṛṣṇā}, \textquotedblleft{}thirst\textquotedblright{} (CDIAL 5941). The same root gives English \emph{thirst}.
\end{cousinweb}

\begin{grammarlens}[title={What you've built}]
One more word for this chapter.
\end{grammarlens}

\subsection*{Your turn}
\begin{itemize}
\raggedright
  \item \emph{Say it:} \emph{tahān}
  \item \emph{Say it:} \emph{tahān}, once more
  \item \emph{Say it:} say \emph{tahān}
  \item \emph{Recall:} say the Marathi for ill, then the Marathi for pain, then say \emph{tahān} again
\end{itemize}

\begin{tcolorbox}[breakable,colback=teal!4,colframe=teal!35!black,title={Before you move on}]
What does \mr{तहान} mean? (\textquotedblleft{}Thirst\textquotedblright{}.) Say it once more.
\end{tcolorbox}
\section[\texorpdfstring{\mr{भूक} (bhūk)}{भूक (bhūk)}]{\mr{भूक} (bhūk) — hunger}

\label{lesson:MR-C73-bhuk}



\begin{quote}
Before the new one: say the Marathi for pain, then the Marathi for thirst.
\end{quote}

\subsection*{You'll want to know: \mr{भूक}}
\textbf{\mr{भूक}} — \emph{bhūk} — \textquotedblleft{}hunger\textquotedblright{}.

\begin{cousinweb}
\textbf{\mr{भूक}} (\emph{bhūk}) is a tadbhava, a word inherited from Sanskrit and reshaped by sound change. It comes from \emph{bubhukṣā}, \textquotedblleft{}the wish to eat, hunger\textquotedblright{} (CDIAL 9286). It is the \textquotedblleft{}wanting\textquotedblright{} form of \emph{bhuj}, \textquotedblleft{}to eat\textquotedblright{}.
\end{cousinweb}

\begin{grammarlens}[title={What you've built}]
One more word for this chapter.
\end{grammarlens}

\subsection*{Your turn}
\begin{itemize}
\raggedright
  \item \emph{Say it:} \emph{bhūk}
  \item \emph{Say it:} \emph{bhūk}, once more
  \item \emph{Say it:} say \emph{bhūk}
  \item \emph{Recall:} say the Marathi for pain, then the Marathi for thirst, then say \emph{bhūk} again
\end{itemize}

\begin{tcolorbox}[breakable,colback=teal!4,colframe=teal!35!black,title={Before you move on}]
What does \mr{भूक} mean? (\textquotedblleft{}Hunger\textquotedblright{}.) Say it once more.
\end{tcolorbox}
\section[\texorpdfstring{\mr{आनंदी} (ānandī)}{आनंदी (ānandī)}]{\mr{आनंदी} (ānandī) — happy}

\label{lesson:MR-C73-anandi}



\begin{quote}
Before the new one: say the Marathi for thirst, then the Marathi for hunger.
\end{quote}

\subsection*{You'll want to know: \mr{आनंदी}}
\textbf{\mr{आनंदी}} — \emph{ānandī} — \textquotedblleft{}happy\textquotedblright{}.

\begin{cousinweb}
\textbf{\mr{आनंदी}} (\emph{ānandī}) is Sanskrit \emph{ānandin}, \textquotedblleft{}happy\textquotedblright{}, a tatsama built on \emph{ānanda}, \textquotedblleft{}happiness\textquotedblright{} (CDIAL 1172).
\end{cousinweb}

\begin{grammarlens}[title={What you've built}]
One more word for this chapter.
\end{grammarlens}

\subsection*{Your turn}
\begin{itemize}
\raggedright
  \item \emph{Say it:} \emph{ānandī}
  \item \emph{Say it:} \emph{ānandī}, once more
  \item \emph{Say it:} say \emph{ānandī}
  \item \emph{Recall:} say the Marathi for thirst, then the Marathi for hunger, then say \emph{ānandī} again
\end{itemize}

\begin{tcolorbox}[breakable,colback=teal!4,colframe=teal!35!black,title={Before you move on}]
What does \mr{आनंदी} mean? (\textquotedblleft{}Happy\textquotedblright{}.) Say it once more.
\end{tcolorbox}
